\section{From Pixel Area to Particle Masses: The Complete OPH Spectrum Derivation}\label{from-pixel-area-to-particle-masses-the-complete-oph-spectrum-derivation}

\begin{quote}
\textbf{Integrated derivation within this manuscript}: this part applies the OPH framework developed in Part I.

\textbf{Abstract.} The Observer-Patch Holography (OPH) framework derives the Standard Model particle spectrum from a single geometric quantity: the \textbf{pixel area} , the area of one computational cell on the holographic screen, measured in Planck units. Starting from the OPH axioms (A1--A4: entanglement equilibrium, MaxEnt, edge-center completion, refinement stability), the ``screen'' encodes all physics through its area law (§5 of Part I of this manuscript) and edge-sector structure (§6 of Part I of this manuscript). The pixel area constant \(P \equiv a_{\text{cell}}/\ell_P^2 = 1.63094\) is \textbf{calibrated} from the gauge coupling sector via the screen\textquotesingle s entropy-matching condition (§5.4 of Part I of this manuscript); a second input, the total screen capacity \(\log(\dim\mathcal{H}) \sim 10^{122}\), enters only for neutrino masses and the cosmological constant (§14 of Part V of this manuscript). Every other quantity appearing in the derivation chain is either a mathematical constant, a consequence of the gauge group (itself reconstructed from the screen\textquotesingle s edge-sector fusion rules via Tannaka-Krein, §6.1 of Part I of this manuscript), or a derived structural integer (\(N_c = 3\), \(N_g = 3\), \(\varepsilon = 1/6\), \(\delta = 2/9\)). Outputs are classified by epistemic status: quantities flowing from the calibrated gauge couplings are consistency checks; quantities determined by independent structural conditions (critical surface, Koide structure, \(\mathbb{Z}_6\) texture) are genuine predictions. A complete audit of all constants and the full epistemic classification appear in §1A--1B below.
\end{quote}

\begin{center}\rule{0.5\linewidth}{0.5pt}\end{center}

\subsection{1. Inputs and Contract}\label{1-inputs-and-contract}

\subsubsection{1.1 Physical Inputs}\label{11-physical-inputs}

The entire derivation pipeline uses exactly \textbf{two} physical inputs:

{\def\LTcaptype{none} % do not increment counter
\begin{longtable}[]{@{}>{\RaggedRight\arraybackslash}p{0.240\textwidth}>{\RaggedRight\arraybackslash}p{0.240\textwidth}>{\RaggedRight\arraybackslash}p{0.240\textwidth}>{\RaggedRight\arraybackslash}p{0.240\textwidth}@{}}
\toprule\noalign{}
Input & Symbol & Value & Origin \\
\midrule\noalign{}
\endhead
\bottomrule\noalign{}
\endlastfoot
Pixel area constant & \(P \equiv a_{\text{cell}}/\ell_P^2\) & 1.63094 & \textbf{Calibrated} from gauge couplings via entropy matching (§5.4 of Part I; see §1B) \\
Screen capacity & \(\log(\dim\mathcal{H}_{\text{tot}})\) & \(\sim 10^{122}\) & Inferred from observed cosmological constant \\
\end{longtable}
}

\subsubsection{1.2 What Counts as "Derived"}\label{12-what-counts-as-derived}

Everything else entering the computation falls into one of three categories:

\begin{enumerate}
\def\labelenumi{\arabic{enumi}.}
\item
  \textbf{Mathematical constants}: \(\pi\), \(e\), \(\zeta(3)\), Casimir eigenvalues, group dimensions, \(\beta\)-function coefficients , all determined by the gauge group structure (itself derived from the axioms via Tannaka-Krein reconstruction).
\item
  \textbf{Derived quantities}: the unification scale \(M_U\), the unified coupling \(\alpha_U\), the electroweak VEV \(v\), all gauge couplings at \(M_Z\), the \(\mathbb{Z}_6\) defect parameter \(\varepsilon = 1/6\), the Koide phase \(\delta = 2/9\), and the Froggatt-Nielsen integer exponents.
\item
  \textbf{Numerical-method parameters}: loop order for RG running (1-loop SM for critical surface, 4-loop MSbar for \(\Lambda_{\overline{\text{MS}}}\)), lattice sizes and statistics for hadron ratios. These affect precision but not the prediction logic.
\end{enumerate}

\subsubsection{1.3 No-Leakage Guarantee}\label{13-no-leakage-guarantee}

The computation code enforces a runtime mutation test: after computing all outputs, the PDG reference table is scrambled and outputs are recomputed. Any change triggers an assertion failure. This verifies that no PDG values leak into the computation pipeline beyond the initial \(P\) calibration. Implemented in \texttt{oph\_predict\_compare.py::\_assert\_pdg\_not\_used()}.

\begin{center}\rule{0.5\linewidth}{0.5pt}\end{center}

\subsection{1A. Complete Audit of All Constants}\label{1a-complete-audit-of-all-constants}

\textbf{Purpose.} This section catalogs \emph{every} named constant, integer, or coefficient appearing anywhere in the prediction code or derivation chain, and explains its origin. The goal is to demonstrate that nothing has been "smuggled in" , every value is either (I) the single physical input \(P\), (II) derived from the OPH axioms under minimal assumptions, (III) a mathematical/group-theoretic constant uniquely fixed by the gauge group, or (IV) a standard QFT result that any textbook reproduces from the Lagrangian.

\subsubsection{1A.1 The Physical Input}\label{1a1-the-physical-input}

{\def\LTcaptype{none} % do not increment counter
\setlength{\LTleft}{0pt}
\setlength{\LTright}{0pt}
\setlength{\tabcolsep}{1.2pt}
\renewcommand{\arraystretch}{1.05}
\footnotesize
\begin{longtable}[]{@{}>{\RaggedRight\arraybackslash}p{0.13\textwidth}>{\RaggedRight\arraybackslash}p{0.09\textwidth}>{\RaggedRight\arraybackslash}p{0.18\textwidth}>{\RaggedRight\arraybackslash}p{0.33\textwidth}>{\RaggedRight\arraybackslash}p{0.15\textwidth}@{}}
\toprule\noalign{}
Constant & Value & Code variable & Origin & Status \\
\midrule\noalign{}
\endhead
\bottomrule\noalign{}
\endlastfoot
\(P \equiv a_{\text{cell}}/\ell_P^2\) & 1.63094 & \texttt{P\_DEFAULT} & The pixel area: area of one UV cell on the holographic screen in Planck units. Extracted from the entropy-matching condition \(P/4 = \bar{\ell}_{\text{SU(2)}}(t_2) + \bar{\ell}_{\text{SU(3)}}(t_3)\) (§5.4 of Part I of this manuscript). This relation is \emph{derived} from the axioms; the \emph{numerical value} is fixed by requiring consistency with the observed gauge couplings (just as Newton\textquotesingle s constant \(G\) is the one free parameter of general relativity). & \textbf{Single free parameter} for particle physics \\
\(\log(\dim\mathcal{H}_{\text{tot}})\) & \(\sim 10^{122}\) & \texttt{LOG\_DIM\_H\_DEFAULT} & Total screen capacity = de Sitter entropy \(S_{dS} = 3\pi/(G\Lambda)\). Enters only for neutrino masses and cosmological constant (§14 of Part V of this manuscript). Derived from the screen area law (§5 of Part I of this manuscript). & \textbf{Second input} (cosmology sector only) \\
\end{longtable}
}

\subsubsection{1A.2 Structural Integers Derived from Axioms}\label{1a2-structural-integers-derived-from-axioms}

{\def\LTcaptype{none} % do not increment counter
\setlength{\LTleft}{0pt}
\setlength{\LTright}{0pt}
\setlength{\tabcolsep}{1.2pt}
\renewcommand{\arraystretch}{1.05}
\footnotesize
\begin{longtable}[]{@{}>{\RaggedRight\arraybackslash}p{0.13\textwidth}>{\RaggedRight\arraybackslash}p{0.12\textwidth}>{\RaggedRight\arraybackslash}p{0.21\textwidth}>{\RaggedRight\arraybackslash}p{0.28\textwidth}>{\RaggedRight\arraybackslash}p{0.14\textwidth}@{}}
\toprule\noalign{}
Constant & Value & Code variable & Derivation & Part I of this manuscript ref \\
\midrule\noalign{}
\endhead
\bottomrule\noalign{}
\endlastfoot
\(N_c\) (colors) & 3 & \texttt{N\_c\_DEFAULT} & Witten\textquotesingle s global SU(2) anomaly requires \(N_c + 1\) SU(2) doublets per generation to be even, so \(N_c\) must be odd. \(N_c = 1\) is excluded (SU(1) trivial). The Selection Axiom MAR (third component of the complexity vector) picks the smallest nontrivial value. \textbf{Result}: \(N_c = 3\). & Theorem 6.14 (§6.9); Part II of this manuscript §6 \\
\(N_g\) (generations) & 3 & \texttt{N\_g\_DEFAULT} & Three independent constraints narrow the window: (1) CP violation requires \(N_g \geq 3\) (CKM phases \((N_g-1)(N_g-2)/2 > 0\)). (2) Asymptotic freedom of SU(2) requires \(N_g(N_c+1) < 22\), giving \(N_g \leq 5\). (3) MAR (fourth component of the complexity vector) picks the smallest value in \(\{3, 4, 5\}\). \textbf{Result}: \(N_g = 3\). & Proposition 6.9 (§6.4); Part II of this manuscript §7 \\
Gauge group & \((SU(3)\times SU(2)\times U(1))/\mathbb{Z}_6\) & --- & Tannaka-Krein / Doplicher-Roberts reconstruction from edge-sector fusion rules yields a compact gauge group \(G\) (Theorem 6.1, §6.1). The Selection Axiom MAR uniquely selects the SM factors: the minimal coupled carrier \(\mathbb{C}^3 \otimes \mathbb{C}^2\) enforces the product structure SU(3) \(\times\) SU(2) \(\times\) U(1), the commutant argument excludes extra factors, and the \(\mathbb{Z}_6\) quotient follows from hypercharge quantization (Proposition 6.6). & §6.1--6.2; Part II of this manuscript \\
\(\beta_{\text{EW}}\) & 4 & \texttt{beta\_ew(N\_c)} & Number of SU(2) doublets per generation = \(N_c\) quark doublets + 1 lepton doublet = \(N_c + 1 = 4\). This is a counting consequence of the Witten anomaly analysis that already fixed \(N_c = 3\). & §6.9, §6.19 \\
\(\varepsilon\) (Z₆ defect) & \(1/6\) & \texttt{defect\_\allowbreak epsilon\_\allowbreak Z6()} & The SM gauge group has a \(\mathbb{Z}_6\) center quotient. Each unit of \(\mathbb{Z}_6\) defect insertion removes \(\Delta S = \ln 6\) nats of entropy (MaxEnt weighting, Assumption B). The resulting Boltzmann suppression per defect is \(e^{-\ln 6} = 1/6\). \textbf{This is topologically fixed} by the quotient structure, not chosen. & §6.18, §6.21 \\
\(\delta\) (Koide phase) & \(2/9\) & computed inline & The holonomy phase on generation space: \(\delta = \beta_{\text{EW}} \cdot Y_Q / N_g = (N_c+1)/(2 N_c N_g) = 4/18 = 2/9\). Here \(Y_Q = 1/(2N_c)\) is the quark-doublet hypercharge (fixed by anomaly cancellation) and \(N_g = 3\) is the number of generations. All ingredients are previously derived. Experimental extraction: \(\delta_{\text{exp}} = 0.2222248 \pm 0.0000063\), matching \(2/9 = 0.2222\ldots\) within \(0.4\sigma\). & Proposition 13.3 (§13 of Part V of this manuscript) \\
Quark exponents & \(n_u = (6,3,0)\), \(n_d = (6,4,2)\) & \texttt{derive\_\allowbreak integer\_\allowbreak vectors()} & Minimal SU(3) hierarchy structure: \(n_{u} = (2N_c, N_c, 0)\) and \(n_d = (2N_c, N_c+1, N_c-1)\). These are the unique sequences that (a) span the range \([0, 2N_c]\) in \(N_g\) steps, (b) have the top quark unsuppressed (\(n_t = 0\)), and (c) match the observed up/down mass ordering under CKM rotation. With \(N_c = 3\): \((6,3,0)\) and \((6,4,2)\). & §6.21 \\
Lepton exponents & \(n_e = (7,4,3)\) & \texttt{derive\_\allowbreak lepton\_\allowbreak exponents()} & Derived algorithmically: the Koide roots \(r_k^2\) (from \(\delta = 2/9\)) must satisfy \(r_i^2/r_j^2 \approx \varepsilon^{n_i - n_j}\) for all pairs. The code scans integer triples near \((N_g, N_g+1, \ldots)\) and selects the unique triple minimizing the log-ratio residuals. \textbf{Result}: \((n_\tau, n_\mu, n_e) = (3, 4, 7)\). & §6.21, §13 of Part V of this manuscript \\
CKM angles & \(s_{12} = \varepsilon\), \(s_{23} = \varepsilon^2\), \(s_{13} = \varepsilon^3\) & inline & Powers of \(\varepsilon = 1/6\): the CKM mixing angles scale as \(\varepsilon^{\lvert n_i - n_j \rvert}\) where the exponents are the generation-gap between up-type and down-type defect charges. This is the standard Froggatt-Nielsen scaling. & §6.21 \\
\(Q\) (Koide ratio) & \(2/3\) & inline & The \(\mathbb{Z}_3\) mode balance on generation space: the circulant Hermitian matrix \(\Phi = a I + b P + b^* P^2\) on 3 generations satisfies \(Q = (1 + 2\lvert b/a \rvert^2)/3 = 2/3\) when \(\lvert b/a \rvert = 1/\sqrt{2}\), i.e., when the singlet and charged \(\mathbb{Z}_3\) modes have equal norm. This is the MaxEnt equilibrium. & Theorem 13.2 (§13 of Part V of this manuscript) \\
\end{longtable}
}

\subsubsection{1A.3 Standard QFT Constants (Fixed by the Gauge Group + Matter Content)}\label{1a3-standard-qft-constants-fixed-by-the-gauge-group--matter-content}

These constants are \emph{not} inputs , they are uniquely determined mathematical consequences of the gauge group SU(3) \(\times\) SU(2) \(\times\) U(1) with \(N_g\) generations of the standard fermion representations. Any QFT textbook (e.g., Peskin \& Schroeder) derives them from the Lagrangian.

{\def\LTcaptype{none} % do not increment counter
\setlength{\LTleft}{0pt}
\setlength{\LTright}{0pt}
\setlength{\tabcolsep}{1.2pt}
\renewcommand{\arraystretch}{1.05}
\footnotesize
\begin{longtable}[]{@{}>{\RaggedRight\arraybackslash}p{0.20\textwidth}>{\RaggedRight\arraybackslash}p{0.16\textwidth}>{\RaggedRight\arraybackslash}p{0.20\textwidth}>{\RaggedRight\arraybackslash}p{0.32\textwidth}@{}}
\toprule\noalign{}
Constant & Value & Code location & What determines it \\
\midrule\noalign{}
\endhead
\bottomrule\noalign{}
\endlastfoot
MSSM \(\beta\)-coefficients & \((b_1, b_2, b_3) = (33/5, 1, -3)\) & \texttt{B\_MSSM} & The edge-sector computation (§6.17 of Part I of this manuscript) derives \(\beta\)-function shifts \(\Delta b \approx (2.49, 4.17, 4.01)\) from the heat-kernel vacuum-polarization weighting and the \(\mathbb{Z}_6\) quotient. These match the shifts from SM to MSSM to better than 1\%. The coefficients themselves are standard 1-loop group theory: \(b_i = \sum_R T(R_i)\) summed over all matter multiplets. Using MSSM content is a \emph{consequence} of the edge-sector matching, not an assumption. SM-only coefficients fail catastrophically (\(\alpha_s\) off by 52\(\sigma\)). \\
SM 1-loop \(\beta\)-coefficients & standard SM set & \texttt{B\_SM} & Standard 1-loop coefficients for the SM with \(N_g = 3\), namely \((41/10,\,-19/6,\,-7)\). Used only for critical-surface RG evolution from \(M_U\) to \(m_t\) (§5 of this paper). Fully fixed by SM gauge group and matter content. \\
4-loop MSbar \(\beta\)-coefficients & \(\beta_0, \beta_1, \beta_2, \beta_3\) & \texttt{beta\_\allowbreak coeffs\_\allowbreak msbar(n\_\allowbreak f)} in \texttt{oph\_\allowbreak qcd.py} & Standard perturbative QCD: \(\beta_0 = 11 - 2n_f/3\), \(\beta_1 = 102 - 38n_f/3\), etc. These are computed from SU(3) Feynman diagrams at each loop order. The 4-loop coefficient \(\beta_3\) involves \(\zeta(3)\); all are universal in the MSbar scheme. \\
Pole-mass conversion & \(K_2 = 13.44 - 1.04 n_l\), \(K_3 = 190.6 - 26.7 n_l + 0.65 n_l^2\) & \texttt{top\_\allowbreak pole\_\allowbreak from\_\allowbreak msbar()} & Standard 3-loop QCD relation between \(\overline{\text{MS}}\) and pole mass (Chetyrkin, Kniehl, Steinhauser 2000). These are perturbative coefficients computed from Feynman diagrams; they depend only on \(n_l\) (number of light flavors) and SU(3) group factors. \\
Casimir eigenvalues & \(C_2(p,q)\), \(d_{(p,q)}\) & \texttt{ellbar\_\allowbreak su3()}, \texttt{ellbar\_\allowbreak su2()} & Group-theoretic: \(C_2(p,q) = \frac{1}{3}(p^2 + q^2 + pq + 3p + 3q)\) for SU(3), \(C_2(j) = j(j+1)\) for SU(2), \(d_{(p,q)} = \frac{1}{2}(p+1)(q+1)(p+q+2)\). These are properties of the Lie algebra, not physical inputs. \\
\(\zeta(3)\) & 1.20206... & \texttt{z3} in \texttt{oph\_\allowbreak qcd.py} & Ap\'{e}ry\textquotesingle s constant , a pure mathematical constant appearing in 4-loop \(\beta_3\). \\
GUT normalization & \(\alpha_1 = \frac{5}{3}\alpha_Y\) & inline & The factor \(5/3\) is the standard GUT normalization ensuring all three couplings unify. It follows from embedding U(1)\(_Y\) into SU(5) and is fixed by the hypercharge assignments of the SM fermions. \\
\end{longtable}
}

\subsubsection{1A.4 Dimensional/Definitional Constants}\label{1a4-dimensionaldefinitional-constants}

{\def\LTcaptype{none} % do not increment counter
\setlength{\LTleft}{0pt}
\setlength{\LTright}{0pt}
\setlength{\tabcolsep}{1.2pt}
\renewcommand{\arraystretch}{1.05}
\footnotesize
\begin{longtable}[]{@{}>{\RaggedRight\arraybackslash}p{0.20\textwidth}>{\RaggedRight\arraybackslash}p{0.16\textwidth}>{\RaggedRight\arraybackslash}p{0.20\textwidth}>{\RaggedRight\arraybackslash}p{0.32\textwidth}@{}}
\toprule\noalign{}
Constant & Value & Code variable & Status \\
\midrule\noalign{}
\endhead
\bottomrule\noalign{}
\endlastfoot
\(E_P\) (Planck energy) & \(1.220890 \times 10^{19}\) GeV & \texttt{E\_PLANCK\_GEV} & \textbf{Definition}: \(E_P = \sqrt{\hbar c^5/G}\). This is a unit conversion factor, not a physical input , it converts from natural (Planck) units to GeV. In the OPH framework, \(G\) itself is derived from the pixel area via \(G = a_{\text{cell}}/(4\bar{\ell}\,\hbar/c^3)\) (§5.4 of Part I of this manuscript), so \(E_P\) is ultimately set by \(P\). \\
\(e^{2\pi}\) in \(M_U\) formula & 535.49... & inline & The Euclidean regularity condition on the collar geometry fixes the angular period to \(2\pi\) (§5.8 of Part I of this manuscript). This is a geometric constant, not a parameter. \\
\(\sqrt{2}\) in Yukawa & \(v/\sqrt{2}\) & inline & Standard Higgs mechanism convention: the Yukawa coupling \(y_f\) relates to the fermion mass via \(m_f = y_f v / \sqrt{2}\). This is a normalization convention, not physics. \\
\(\Delta\rho_{\text{stage-3}}\) & \(3/(32\pi^2) \approx 0.0095\) & inline & Universal 1-loop custodial correction from a unit-Yukawa fermion doublet. This is a standard EW radiative correction; it depends only on the gauge structure, not on measured masses. \\
\end{longtable}
}

\subsubsection{1A.5 Numerical-Method Parameters}\label{1a5-numerical-method-parameters}

These affect computational precision but not the prediction logic. Changing them changes the last digits, not the physics.

{\def\LTcaptype{none} % do not increment counter
\begin{longtable}[]{@{}>{\RaggedRight\arraybackslash}p{0.320\textwidth}>{\RaggedRight\arraybackslash}p{0.320\textwidth}>{\RaggedRight\arraybackslash}p{0.320\textwidth}@{}}
\toprule\noalign{}
Parameter & Value used & Effect of changing \\
\midrule\noalign{}
\endhead
\bottomrule\noalign{}
\endlastfoot
Loop order for RG (critical surface) & 1-loop SM & 2-loop would shift \(m_H\) by \textasciitilde1 GeV, \(m_t\) by \textasciitilde0.5 GeV \\
Loop order for \(\Lambda_{\overline{\text{MS}}}\) & 4-loop MSbar & 3-loop changes \(\Lambda\) by \textasciitilde2\% \\
RK4 step count & 2000 steps & Doubling changes nothing to displayed precision \\
Lattice sizes (hadrons) & \(L = 2\)--\(6\) & Larger volumes needed for precision; current values are prototype \\
Bisection tolerance & \(10^{-10}\) & Affects last digits of \(\alpha_U\) only \\
Heat-kernel truncation & 200 representations & Adding more changes \(\bar{\ell}\) by \(< 10^{-12}\) \\
\end{longtable}
}

\subsubsection{1A.6 Summary: What Is and Isn\textquotesingle t Assumed}\label{1a6-summary-what-is-and-isnt-assumed}

\textbf{Derived from OPH axioms (no assumptions beyond A1--A4)}:

\begin{itemize}
\tightlist
\item
  Gauge group structure (Tannaka-Krein reconstruction from edge sectors)
\item
  \(\varepsilon = 1/6\) (topological, from \(\mathbb{Z}_6\) quotient)
\item
  \(\lambda(M_U) = 0\), \(\beta_\lambda(M_U) = 0\) (refinement stability)
\item
  \(Q = 2/3\) (MaxEnt on \(\mathbb{Z}_3\) generation space)
\item
  Heat-kernel form \(p_R \propto d_R e^{-tC_2(R)}\) (MaxEnt + edge completion)
\end{itemize}

\textbf{Derived under the Selection Axiom MAR (Part II of this manuscript)}:

\begin{itemize}
\tightlist
\item
  SM gauge group \(SU(3) \times SU(2) \times U(1)/\mathbb{Z}_6\) (MAR + admissibility)
\item
  \(N_c = 3\) (Witten anomaly + MAR minimality)
\item
  \(N_g = 3\) (CP violation + asymptotic freedom + MAR minimality)
\item
  \(\delta = 2/9\) (algebraic consequence of \(N_c = 3\), \(N_g = 3\), \(\beta_{\text{EW}} = 4\))
\item
  \(\beta_{\text{EW}} = 4\) (doublet counting from \(N_c = 3\))
\item
  Integer exponents \(n_u\), \(n_d\), \(n_e\) (algorithmically from \(N_c\), \(N_g\), \(\varepsilon\), \(\delta\))
\end{itemize}

\textbf{Fixed by the gauge group (standard QFT, no free parameters)}:

\begin{itemize}
\tightlist
\item
  All \(\beta\)-function coefficients (1-loop through 4-loop)
\item
  Pole-mass conversion coefficients \(K_2\), \(K_3\)
\item
  Casimir eigenvalues, representation dimensions
\item
  GUT normalization factor \(5/3\)
\end{itemize}

\textbf{The single free parameter}:

\begin{itemize}
\tightlist
\item
  \(P = 1.63094\) (pixel area), calibrated from the gauge coupling sector (see §1B)
\end{itemize}

\textbf{Second input (cosmology only)}:

\begin{itemize}
\tightlist
\item
  \(\log(\dim\mathcal{H}) \sim 10^{122}\) (screen capacity)
\end{itemize}

\textbf{What is NOT assumed or smuggled in}:

\begin{itemize}
\tightlist
\item
  No particle masses
\item
  No CKM matrix elements
\item
  No Higgs VEV or quartic coupling
\item
  No SUSY-breaking scale
\item
  No Yukawa couplings
\item
  No \(\Lambda_{\text{QCD}}\)
\end{itemize}

\begin{center}\rule{0.5\linewidth}{0.5pt}\end{center}

\subsection{1B. Calibration vs.\ Prediction: Epistemic Classification of Outputs}\label{1b-calibration-vs-prediction}

\subsubsection{1B.1 The Calibration Structure}\label{1b1-the-calibration-structure}

The pixel area constant \(P\) is calibrated from the entropy-matching condition \(P/4 = \bar{\ell}_{\text{SU(2)}}(t_2) + \bar{\ell}_{\text{SU(3)}}(t_3)\), where \(t_i = 4\pi^2 \alpha_i\) depends on the gauge couplings. This means \(P\) encodes information from the gauge coupling sector. Any quantity that flows directly from the gauge couplings through the same algebraic chain used to calibrate \(P\) is therefore a \textbf{consistency check}---it must reproduce the input data by construction, and its agreement with experiment does not constitute an independent test of the framework.

This is not a defect of OPH any more than \(G\) being calibrated from gravitational measurements is a defect of general relativity. Every quantitative theory requires at least one calibration to set its scale. The question is: how many independent outputs does the framework produce beyond the calibration sector?

\subsubsection{1B.2 Why ``Downstream of \(P\)'' Does Not Mean ``Determined by \(P\)''}\label{1b2-dimensionless-independence}

A common objection is that every quantity computed using \(P\) ``inherits the calibration data,'' rendering all outputs circular. This conflates two distinct roles that \(P\) plays:

\begin{enumerate}
\item \textbf{Scale setting.} \(P\) determines the overall energy scale (via \(\alpha_U \to v\)). Every theory has this: \(G\) sets the Planck scale in general relativity, \(\alpha\) sets the Bohr radius in QED. Quantities that depend on \(P\) only for their conversion from dimensionless ratios to GeV are not ``fitted''---they are \emph{scaled}, exactly as a QED cross-section is scaled by \(\alpha\) without being circular.
\item \textbf{Structural determination.} The calibration procedure solves \emph{one} equation (the pixel constraint) for \emph{one} unknown (\(\alpha_U\)). This single degree of freedom determines the gauge couplings, \(v\), \(M_U\), \(m_W\), and \(m_Z\)---all members of a one-parameter family. These are the consistency checks (Tier~1).
\end{enumerate}

The Tier~2 predictions are determined by structural content---boundary conditions, algebraic identities, topological data---that is \emph{logically independent} of the pixel constraint. Their dimensionless values are fixed before \(P\) enters the computation. Concretely:

{\footnotesize
\def\LTcaptype{none}
\begin{longtable}[]{@{}>{\RaggedRight\arraybackslash}p{0.22\textwidth}>{\RaggedRight\arraybackslash}p{0.30\textwidth}>{\RaggedRight\arraybackslash}p{0.40\textwidth}@{}}
\toprule
Prediction & Dimensionless form & What determines it (not \(P\)) \\
\midrule
\endhead
\bottomrule
\endfoot
\(m_e : m_\mu : m_\tau\) & Koide ratios from \(Q = 2/3\), \(\delta = 2/9\) & Algebraic: \(\delta = (N_c+1)/(2N_cN_g)\), \(Q\) from MaxEnt on \(\mathbb{Z}_3\). Zero free parameters. \\
\(m_H / v\) & \(\approx 0.509\); from \(\lambda(M_U) = 0\), \(\beta_\lambda(M_U) = 0\) & Critical surface (refinement stability). Independent boundary condition, not used in calibrating \(P\). \\
\(m_t / v\) & \(\approx 0.702\); from \(y_t(M_U)\) at \(\beta_\lambda = 0\) & Same critical surface. \\
\(m_f / m_t\) (quarks) & \(\varepsilon^{n_f}\) with \(\varepsilon = 1/6\), integer \(n_f\) & \(\mathbb{Z}_6\) quotient topology + \(N_c\), \(N_g\). No continuous parameter. \\
\end{longtable}
}

\textbf{The decisive test}: change \(P\) to any other value. The lepton mass \emph{ratios} do not change. The Higgs-to-VEV ratio does not change. The quark exponents do not change. Only the overall energy scale (the value of \(v\) in GeV) shifts. This is the operational definition of ``independent of the calibration.''

To see why, note that \(P\) is a single real number. It cannot simultaneously encode the electron-to-muon mass ratio (\(m_e/m_\mu \approx 1/206.8\)), the Higgs-to-top ratio (\(m_H/m_t \approx 0.725\)), and the six quark-mass exponents \((6,3,0,6,4,2)\). These are independent degrees of freedom whose values are fixed by structural constraints that do not reference \(P\), the entropy-matching condition, or any experimental coupling constant. Their agreement with observation is a test the framework could fail.

\subsubsection{1B.3 Classification}\label{1b3-classification}

\textbf{Tier 1---Calibration-sector consistency checks.} These quantities are algebraically entangled with \(P\)'s calibration. Their agreement with PDG data demonstrates internal consistency but does not constitute independent prediction.

{\footnotesize
\begin{longtable}[]{@{}>{\RaggedRight\arraybackslash}p{0.30\textwidth}>{\RaggedRight\arraybackslash}p{0.62\textwidth}@{}}
\toprule
Quantity & Why it is a consistency check \\
\midrule
\endhead
\bottomrule
\endfoot
\(\alpha_s(m_Z)\), \(\alpha_{\text{em}}^{-1}(m_Z)\), \(\sin^2\theta_W(m_Z)\) & The gauge couplings at \(m_Z\) are the data from which \(P\) (via \(\alpha_U\)) is extracted. The self-consistent fixed-point procedure ensures these reproduce. \\
\(m_Z\) (tree-level) & \(m_Z = \frac{1}{2}v\sqrt{g_2^2 + g_Y^2}\) where \(v\) and the couplings both flow from \(\alpha_U\), determined by the pixel constraint that calibrated \(P\). \\
\(m_W\) (tree-level) & Same: \(m_W = \frac{1}{2}v \cdot g_2\), directly from the calibrated couplings. \\
\(v\) (Higgs VEV) & \(v = E_{\text{cell}} \cdot e^{-2\pi/(\beta_{\text{EW}} \alpha_U)}\); determined by \(\alpha_U\), the variable solved for during \(P\) calibration. \\
\end{longtable}
}

\textbf{Tier 2---Genuine predictions.} These quantities are determined by structural content of OPH (critical surface conditions, Koide structure, \(\mathbb{Z}_6\) texture) that is logically independent of the gauge-coupling calibration. They use \(P\) only to set an overall scale, and their specific values are not guaranteed by the calibration.

{\footnotesize
\begin{longtable}[]{@{}>{\RaggedRight\arraybackslash}p{0.30\textwidth}>{\RaggedRight\arraybackslash}p{0.62\textwidth}@{}}
\toprule
Quantity & Why it is a genuine prediction \\
\midrule
\endhead
\bottomrule
\endfoot
\(m_H\) (Higgs mass) & From \(\lambda(M_U) = 0\), \(\beta_\lambda(M_U) = 0\) (refinement stability)---an independent boundary condition not used in calibrating \(P\). \\
\(m_t\) (top mass, critical surface) & From \(y_t(M_U)\) fixed by \(\beta_\lambda = 0\); independent of the pixel constraint. \\
\(m_e, m_\mu, m_\tau\) (charged leptons) & From the Koide structure (\(Q = 2/3\), \(\delta = 2/9\))---algebraically independent of the gauge-coupling sector. \\
Quark mass hierarchy & From \(\mathbb{Z}_6\) texture with exponents derived from \(N_c\), \(N_g\); independent of \(P\)'s value. \\
\(m_{\nu_1}, m_{\nu_2}, m_{\nu_3}\) & From screen capacity \(\log(\dim\mathcal{H})\), a second input entirely independent of \(P\). \\
\end{longtable}
}

\textbf{Tier 3---Parameter-free structural predictions.} These discrete quantities follow from the axioms alone with no continuous free parameter.

{\footnotesize
\begin{longtable}[]{@{}>{\RaggedRight\arraybackslash}p{0.45\textwidth}>{\RaggedRight\arraybackslash}p{0.45\textwidth}@{}}
\toprule
Quantity & Origin \\
\midrule
\endhead
\bottomrule
\endfoot
Gauge group \(SU(3) \times SU(2) \times U(1)/\mathbb{Z}_6\) & Tannaka-Krein + MAR \\
\(N_c = 3\), \(N_g = 3\) & Anomaly cancellation + MAR \\
\(\varepsilon = 1/6\) & \(\mathbb{Z}_6\) quotient topology \\
\(\delta = 2/9\), \(Q = 2/3\) & Algebraic from \(N_c\), \(N_g\), \(\beta_{\text{EW}}\) \\
\(m_\gamma = m_g = m_{\text{graviton}} = 0\) & Symmetry protection \\
Proton stability & Product gauge group (no leptoquark generators) \\
\end{longtable}
}

\subsubsection{1B.4 What OPH Achieves}\label{1b4-what-oph-achieves}

The Standard Model has 19 free parameters (or 26 with neutrino masses). OPH reduces this to 1 continuous free parameter (\(P\)) plus 1 cosmological input (\(\log \dim \mathcal{H}\)), with all structural content (gauge group, generation count, mass hierarchy pattern) derived from axioms.

The genuine predictive content---Tier 2---includes the Higgs mass, top mass, all three charged lepton masses, the fermion mass hierarchy, and neutrino masses. The charged leptons achieve sub-permille accuracy from a structure (\(Q = 2/3\), \(\delta = 2/9\)) that has no free parameters. The Higgs and top masses are predicted at the \textasciitilde1\% level from a boundary condition (\(\lambda = \beta_\lambda = 0\) at \(M_U\)) that is independent of the gauge-coupling calibration.

The quark masses are the weakest sector: the hierarchy is qualitatively correct (5 orders of magnitude from integer exponents in base 6), but individual masses show 16--73\% deviations due to missing scheme matching, order-one Clebsch factors, and \(u\)-\(d\) degeneracy.

\(P\) is analogous to Newton's constant \(G\) in general relativity, or \(\alpha\) in QED: a single dimensionful parameter that must be measured. The framework's claim is not ``we predict everything from nothing'' but rather: \textbf{from one measured constant and a set of axioms, the framework determines the complete structure of particle physics, with most outputs being genuine predictions independent of the calibration.}

The gauge couplings at \(m_Z\) are consistency checks. The lepton masses, Higgs mass, top mass, and mass hierarchy are predictions. The quark masses are qualitative.

\begin{center}\rule{0.5\linewidth}{0.5pt}\end{center}

\subsection{2. Stage 1: Fundamental Scales}\label{2-stage-1-fundamental-scales}

\subsubsection{2.1 Unification Scale}\label{21-unification-scale}

The OPH framework derives a unification scale from the pixel area:

\[M_U = \frac{E_P}{e^{2\pi}} \cdot P^{1/6}\]

where \(E_P = 1.220890 \times 10^{19}\) GeV is the Planck energy. This gives:

\[M_U \approx 2.474 \times 10^{16} \text{ GeV}\]

\textbf{Derivation}: The factor \(e^{2\pi}\) arises from the modular geometry of the collar (the Euclidean regularity condition fixes the angular period to \(2\pi\)). The \(P^{1/6}\) factor comes from the dimensional relation between pixel area and the UV cell scale (§5.8 of Part I of this manuscript).

\subsubsection{2.2 Cell Energy Scale}\label{22-cell-energy-scale}

The UV cell energy is:

\[E_{\text{cell}} = \frac{E_P}{\sqrt{P}} \approx 9.56 \times 10^{18} \text{ GeV}\]

This is the natural energy scale of a single computational element on the holographic screen.

\begin{center}\rule{0.5\linewidth}{0.5pt}\end{center}

\subsection{3. Stage 2: Gauge Closure and the Pixel Constraint}\label{3-stage-2-gauge-closure-and-the-pixel-constraint}

\subsubsection{3.1 The Heat-Kernel Entropy}\label{31-the-heat-kernel-entropy}

The edge-center completion (Theorem 2.3 of Part I of this manuscript) combined with MaxEnt yields sector probabilities of the heat-kernel form:

\[p_R(t) \propto d_R \, e^{-t \, C_2(R)}\]

where \(d_R\) is the representation dimension, \(C_2(R)\) the quadratic Casimir, and \(t = 4\pi^2 \alpha\) is the diffusion parameter encoding the gauge coupling.

The \textbf{mean entropy per cell} (the \(\bar{\ell}\) function) for each gauge factor is:

\[\bar{\ell}_G(t) = \sum_R p_R(t) \log d_R\]

For SU(2):

\[\bar{\ell}_{\text{SU(2)}}(t) = \sum_{j=0,1/2,1,...} p_j(t) \log(2j+1), \quad p_j \propto (2j+1) e^{-t \, j(j+1)}\]

For SU(3):

\[\bar{\ell}_{\text{SU(3)}}(t) = \sum_{p,q \geq 0} p_{(p,q)}(t) \log d_{(p,q)}, \quad p_{(p,q)} \propto d_{(p,q)} \, e^{-t \, C_2(p,q)}\]

where \(d_{(p,q)} = \frac{1}{2}(p+1)(q+1)(p+q+2)\) and \(C_2(p,q) = \frac{1}{3}(p^2 + q^2 + pq + 3p + 3q)\).

\subsubsection{3.2 The Pixel Constraint}\label{32-the-pixel-constraint}

The generalized entropy matching (§5.4 of Part I of this manuscript) gives:

\[\frac{P}{4} = \bar{\ell}_{\text{SU(2)}}(t_2) + \bar{\ell}_{\text{SU(3)}}(t_3)\]

where \(t_i = 4\pi^2 \alpha_i\). This is the \textbf{pixel constraint}: the total edge entropy per cell must equal \(P/4\).

\subsubsection{\texorpdfstring{3.3 Gauge Running and the \(\alpha_U\) Solution}{3.3 Gauge Running and the \textbackslash alpha\_U Solution}}\label{33-gauge-running-and-the-alpha_u-solution}

At the unification scale, all three gauge couplings emerge from a single \(\alpha_U\). Running down to a scale \(\mu\) with MSSM-like \(\beta\)-function coefficients \((b_1, b_2, b_3) = (33/5, 1, -3)\):

\[\alpha_i^{-1}(\mu) = \alpha_U^{-1} + \frac{b_i}{2\pi} \ln\frac{M_U}{\mu}\]

\textbf{Why MSSM coefficients?} The edge-sector computation (§6.17 of Part I of this manuscript) derives \(\beta\)-function shifts \(\Delta b \approx (2.49, 4.17, 4.01)\) from the heat-kernel form, the \(\mathbb{Z}_6\) quotient structure, and the Peter-Weyl vacuum-polarization weighting. These match the MSSM shifts \((2.5, 4.17, 4.0)\) to better than 1\%. SM-only coefficients catastrophically fail (\(\alpha_s\) prediction off by 52\(\sigma\)).

\subsubsection{3.4 Self-Consistent Fixed Point}\label{34-self-consistent-fixed-point}

The system is solved self-consistently:

\begin{enumerate}
\def\labelenumi{\arabic{enumi}.}
\tightlist
\item
  \textbf{Trial \(\alpha_U\)} \(\rightarrow\) run couplings to scale \(\mu\).
\item
  \textbf{Compute} \(v = E_{\text{cell}} \cdot \exp(-2\pi/(\beta_{\text{EW}} \cdot \alpha_U))\) with \(\beta_{\text{EW}} = N_c + 1 = 4\).
\item
  \textbf{Compute} tree-level \(m_Z(\mu) = \frac{1}{2} v \sqrt{g_2^2 + g_Y^2}\) where \(g_Y = \sqrt{4\pi \cdot \frac{3}{5}\alpha_1}\).
\item
  \textbf{Find} \(\mu^* = m_Z(\mu^*)\) (self-consistent fixed point).
\item
  \textbf{Evaluate} the pixel residual: \(\bar{\ell}_{\text{SU(2)}}(t_2) + \bar{\ell}_{\text{SU(3)}}(t_3) - P/4\).
\item
  \textbf{Bisect} on \(\alpha_U\) until the pixel residual vanishes.
\end{enumerate}

This yields:

\[\alpha_U \approx 0.04112, \quad \alpha_U^{-1} \approx 24.32\]

\begin{center}\rule{0.5\linewidth}{0.5pt}\end{center}

\subsection{4. Stage 3: Electroweak Observables}\label{4-stage-3-electroweak-observables}

With \(\alpha_U\) and \(M_U\) determined, the gauge couplings at \(\mu = m_{Z,\text{run}}\) are fixed:

{\def\LTcaptype{none} % do not increment counter
\begin{longtable}[]{@{}>{\RaggedRight\arraybackslash}p{0.320\textwidth}>{\RaggedRight\arraybackslash}p{0.320\textwidth}>{\RaggedRight\arraybackslash}p{0.320\textwidth}@{}}
\toprule\noalign{}
Quantity & Formula & Predicted Value \\
\midrule\noalign{}
\endhead
\bottomrule\noalign{}
\endlastfoot
\(m_{Z,\text{run}}\) & Self-consistent fixed point & 91.652 GeV \\
\(v\) (Higgs VEV) & \(E_{\text{cell}} \cdot e^{-2\pi/(\beta_{\text{EW}} \alpha_U)}\) & 246.77 GeV \\
\(\alpha_1(m_Z)\) & \([\alpha_U^{-1} + \frac{b_1}{2\pi}\ln(M_U/m_Z)]^{-1}\) & 0.01696 \\
\(\alpha_2(m_Z)\) & \([\alpha_U^{-1} + \frac{b_2}{2\pi}\ln(M_U/m_Z)]^{-1}\) & 0.03384 \\
\(\alpha_3(m_Z)\) & \([\alpha_U^{-1} + \frac{b_3}{2\pi}\ln(M_U/m_Z)]^{-1}\) & 0.1183 \\
\end{longtable}
}

From these:

\[\alpha_{\text{em}} = \left(\frac{1}{\alpha_2} + \frac{1}{\frac{3}{5}\alpha_1}\right)^{-1}, \quad \sin^2\theta_W = \frac{\alpha_{\text{em}}}{\alpha_2}\]

{\def\LTcaptype{none} % do not increment counter
\begin{longtable}[]{@{}>{\RaggedRight\arraybackslash}p{0.320\textwidth}>{\RaggedRight\arraybackslash}p{0.320\textwidth}>{\RaggedRight\arraybackslash}p{0.320\textwidth}@{}}
\toprule\noalign{}
Observable & Predicted & PDG Value \\
\midrule\noalign{}
\endhead
\bottomrule\noalign{}
\endlastfoot
\(\alpha_{\text{em}}^{-1}(m_Z)\) & 128.31 & 127.952 \\
\(\sin^2\theta_W(m_Z)\) & 0.2307 & 0.23122 \\
\(\alpha_s(m_Z)\) & 0.1183 & 0.1179 \\
\end{longtable}
}

\subsubsection{4.1 Z Boson Pole Mass}\label{41-z-boson-pole-mass}

The tree-level \(m_{Z,\text{run}} = 91.65\) GeV is the running mass at its own scale. The physical pole mass includes the custodial correction from top-quark loops:

\[\Delta\rho_{\text{stage-3}} = \frac{3}{32\pi^2} \approx 0.009499\]

\[M_{Z,\text{pole}} = \frac{m_{Z,\text{run}}}{\sqrt{1 + \Delta\rho}} \approx 91.220 \text{ GeV}\]

\textbf{Why \(\Delta\rho = 3/(32\pi^2)\)?} This is the universal one-loop custodial-symmetry-breaking contribution from a unit Yukawa coupling (\(y_t = 1\)), requiring no knowledge of the actual top mass , only that the top Yukawa is order-one (the least-suppressed channel in the \(\mathbb{Z}_6\) texture).

\subsubsection{4.2 W Boson Mass}\label{42-w-boson-mass}

At tree level:

\[m_W = \frac{1}{2} v \cdot g_2 = \frac{1}{2} v \sqrt{4\pi\alpha_2} \approx 80.39 \text{ GeV}\]

\begin{center}\rule{0.5\linewidth}{0.5pt}\end{center}

\subsection{5. Stage 4: Critical Surface , Higgs and Top}\label{5-stage-4-critical-surface--higgs-and-top}

\subsubsection{5.1 The Critical Surface Constraint}\label{51-the-critical-surface-constraint}

Refinement stability (§6.3, §6.22 of Part I of this manuscript) pushes the Higgs quartic coupling to a marginal stability point at the unification scale:

\[\lambda(M_U) = 0, \quad \beta_\lambda(M_U) = 0\]

Setting \(\beta_\lambda = 0\) at one loop with \(\lambda = 0\) fixes the top Yukawa boundary condition:

\[y_t(M_U) = \left[\frac{1}{16}\left(2g_2^4 + (g_2^2 + g_1^2)^2\right)\right]^{1/4}\]

where \(g_1(M_U)\) and \(g_2(M_U)\) are obtained by running the OPH-derived couplings from \(m_Z\) to \(M_U\) using 1-loop SM \(\beta\)-functions.

\subsubsection{5.2 RG Evolution to Low Scales}\label{52-rg-evolution-to-low-scales}

The coupled system \((y_t, \lambda, g_1, g_2, g_3)\) is integrated from \(M_U\) down to \(\mu \sim m_t\) using RK4 with 1-loop SM \(\beta\)-functions. The gauge couplings run analytically:

\[\alpha_i^{-1}(\mu) = \alpha_i^{-1}(m_Z) - \frac{b_i^{\text{SM}}}{2\pi}\ln(\mu/m_Z)\]

with SM coefficients \((b_1, b_2, b_3)^{\text{SM}} = (41/10, -19/6, -7)\).

\subsubsection{5.3 Top Mass}\label{53-top-mass}

The running top mass is determined self-consistently: iterate \(\mu_{\text{guess}} = y_t(\mu) \cdot v/\sqrt{2}\) until convergence:

\[m_t^{\overline{\text{MS}}}(m_t) \approx 160.6 \text{ GeV}\]

The pole mass includes the standard 3-loop QCD relation:

\[\frac{m_t^{\text{pole}}}{m_t^{\overline{\text{MS}}}} = 1 + \frac{4}{3}\frac{\alpha_s}{\pi} + K_2 \left(\frac{\alpha_s}{\pi}\right)^2 + K_3 \left(\frac{\alpha_s}{\pi}\right)^3\]

with \(K_2 = 13.44 - 1.04 n_l\) and \(K_3 = 190.6 - 26.7 n_l + 0.65 n_l^2\) (\(n_l = 5\)).

\[m_t^{\text{pole}} \approx 171.1 \text{ GeV}\]

\subsubsection{5.4 Higgs Mass}\label{54-higgs-mass}

From \(\lambda(m_t)\) obtained by the RG integration:

\[m_H = \sqrt{2\lambda(m_t)} \cdot v \approx 126.5 \text{ GeV}\]

\begin{center}\rule{0.5\linewidth}{0.5pt}\end{center}

\subsection{6. Stage 5: Discrete Spectrum , Quarks and Leptons}\label{6-stage-5-discrete-spectrum--quarks-and-leptons}

\subsubsection{\texorpdfstring{6.1 The \(\mathbb{Z}_6\) Defect Parameter}{6.1 The \textbackslash mathbb\{Z\}\_6 Defect Parameter}}\label{61-the-mathbbz_6-defect-parameter}

The SM global gauge group is \((SU(3) \times SU(2) \times U(1))/\mathbb{Z}_6\). The entropy deficit from the \(\mathbb{Z}_6\) quotient yields the universal suppression factor:

\[\varepsilon = e^{-\ln 6} = \frac{1}{6}\]

This is the base of the Froggatt-Nielsen texture: each unit of \(\mathbb{Z}_6\) defect insertion suppresses a Yukawa coupling by a factor of 6.

\subsubsection{6.2 Integer Exponents (Algorithmically Derived)}\label{62-integer-exponents-algorithmically-derived}

The diagonal mass exponents are determined by the minimal SU(3) hierarchy structure with \(N_c = 3\) colors and \(N_g = 3\) generations:

\textbf{Up-type quarks}: \(n_u = (2N_c, N_c, 0) = (6, 3, 0)\)

\textbf{Down-type quarks}: \(n_d = (2N_c, N_c+1, N_c-1) = (6, 4, 2)\)

\textbf{Charged leptons}: derived from the Koide phase.

\subsubsection{6.3 The Koide Phase}\label{63-the-koide-phase}

The charged lepton masses use a circulant ansatz on generation space with Koide ratio \(Q = 2/3\) (from \(\mathbb{Z}_3\) mode balance). The holonomy phase is:

\[\delta = \frac{\beta_{\text{EW}}}{2 N_c N_g} = \frac{N_c + 1}{2 N_c N_g} = \frac{4}{18} = \frac{2}{9}\]

The Koide roots are:

\[r_k = 1 + \sqrt{2}\cos\left(\frac{2}{9} + \frac{2\pi k}{3}\right), \quad k = 0, 1, 2\]

sorted in ascending order. The lepton exponents \((n_e, n_\mu, n_\tau)\) are selected by matching ratios \(r_i^2/r_j^2 \approx \varepsilon^{n_i - n_j}\), yielding \((n_e, n_\mu, n_\tau) = (7, 4, 3)\).

\subsubsection{6.4 Diagonal Quark Masses}\label{64-diagonal-quark-masses}

\[m_{u_i} = \frac{v}{\sqrt{2}} \varepsilon^{n_{u,i}}, \quad m_{d_i} = \frac{v}{\sqrt{2}} \varepsilon^{n_{d,i}}\]

\subsubsection{6.5 CKM Mixing}\label{65-ckm-mixing}

The CKM matrix is parameterized by \(\varepsilon\):

\[s_{12} = \varepsilon = 1/6, \quad s_{23} = \varepsilon^2 = 1/36, \quad s_{13} = \varepsilon^3 = 1/216\]

Physical down-type masses are the singular values of \(V_{\text{CKM}} \cdot \text{diag}(m_d, m_s, m_b)\).

\subsubsection{6.6 Charged Lepton Masses}\label{66-charged-lepton-masses}

A scale factor \(S\) is determined from the exponents \(n_e = (7, 4, 3)\) and the Koide roots by demanding internal consistency:

\[\ln S_0 = -\frac{1}{3}\sum_{i=1}^3 \ln\left(\frac{r_i^2 \sqrt{2}}{v} \cdot 6^{n_{e,i}}\right)\]

Then \(S = S_0 \cdot 2^{1/6}\) and:

\[m_e = S \cdot r_1^2, \quad m_\mu = S \cdot r_2^2, \quad m_\tau = S \cdot r_3^2\]

\begin{center}\rule{0.5\linewidth}{0.5pt}\end{center}

\subsection{7. Stage 6: QCD Scale and Hadron Masses}\label{7-stage-6-qcd-scale-and-hadron-masses}

\subsubsection{\texorpdfstring{7.1 \(\Lambda_{\overline{\text{MS}}}\) from \(\alpha_s(m_Z)\)}{7.1 \textbackslash Lambda\_\{\textbackslash overline\{\textbackslash text\{MS\}\}\} from \textbackslash alpha\_s(m\_Z)}}\label{71-lambda_overlinetextms-from-alpha_sm_z}

Given the OPH-derived \(\alpha_s(m_Z) \approx 0.1183\) (a calibration-sector output) and the quark thresholds (\(m_b\), \(m_c\) from Stage 5), the QCD scale is extracted using the standard 4-loop MSbar definition:

\[\ln\frac{\mu^2}{\Lambda^2} = \frac{1}{\beta_0 a} + \frac{\beta_1}{\beta_0^2}\ln(\beta_0 a) + \int_0^a dx\left[\frac{1}{\beta(x)} + \frac{1}{\beta_0 x^2} - \frac{\beta_1}{\beta_0^2 x}\right]\]

where \(a = \alpha_s/(4\pi)\) and \(\beta_0, \beta_1, \beta_2, \beta_3\) are the 4-loop MSbar coefficients.

Stepping down across flavor thresholds:

\[\alpha_s^{(n_f-1)}(\mu_{\text{th}}) = \alpha_s^{(n_f)}(\mu_{\text{th}}) \quad \text{at } \mu_{\text{th}} = m_b, m_c\]

yields:

{\def\LTcaptype{none} % do not increment counter
\begin{longtable}[]{@{}>{\RaggedRight\arraybackslash}p{0.480\textwidth}>{\RaggedRight\arraybackslash}p{0.480\textwidth}@{}}
\toprule\noalign{}
\(n_f\) & \(\Lambda_{\overline{\text{MS}}}^{(n_f)}\) (GeV) \\
\midrule\noalign{}
\endhead
\bottomrule\noalign{}
\endlastfoot
5 & 0.211 \\
4 & 0.288 \\
3 & 0.322 \\
\end{longtable}
}

\subsubsection{7.2 Hadron Masses}\label{72-hadron-masses}

Hadron masses require a nonperturbative computation of dimensionless ratios \(C_X = m_X / \Lambda^{(3)}\). These are computed by an internal lattice QCD collar calculation:

\begin{enumerate}
\def\labelenumi{\arabic{enumi}.}
\tightlist
\item
  \textbf{Wilson SU(3) gauge} (quenched) with Metropolis updates.
\item
  \textbf{Wilson valence quarks} at multiple \(\kappa\) values for chiral extrapolation.
\item
  \textbf{Gradient-flow coupling} for input-free \(a\Lambda\) determination.
\item
  \textbf{Two-point Richardson extrapolation} for the continuum limit.
\end{enumerate}

Then: \(m_X = C_X \cdot \Lambda_{\overline{\text{MS}}}^{(3)}\).

The lattice computation uses small volumes (\(L = 2\)--\(6\)) in the quenched approximation. Systematic uncertainties from quenching (\textasciitilde10\%), finite volume, and limited statistics dominate the error budget for hadron masses.

\begin{center}\rule{0.5\linewidth}{0.5pt}\end{center}

\subsection{8. Stage 7: Neutrino Masses}\label{8-stage-7-neutrino-masses}

\subsubsection{8.1 The Capacity Model}\label{81-the-capacity-model}

The second OPH input , the screen capacity \(\log(\dim\mathcal{H}) \sim 10^{122}\) , enters through the cosmological constant:

\[\rho_\Lambda = \frac{3}{8} \frac{M_P^4}{\log(\dim\mathcal{H})}\]

The neutrino mass scale is anchored at \(\rho_\Lambda^{1/4}\):

\[m_{\nu_3} = \rho_\Lambda^{1/4} \approx 3.0 \times 10^{-12} \text{ GeV} \approx 3.0 \text{ meV}\]

The hierarchy uses the same \(\varepsilon = 1/6\) suppression:

\[m_{\nu_2} = \varepsilon \cdot m_{\nu_3} \approx 0.50 \text{ meV}, \quad m_{\nu_1} = \varepsilon^2 \cdot m_{\nu_3} \approx 0.084 \text{ meV}\]

This yields \(\Delta m_{32}^2 \approx 9.1 \times 10^{-24} \text{ GeV}^2\) and \(\Delta m_{21}^2 \approx 2.5 \times 10^{-25} \text{ GeV}^2\), consistent with the observed atmospheric and solar mass splittings in order of magnitude.

\begin{center}\rule{0.5\linewidth}{0.5pt}\end{center}

\subsection{9. Massless Particles}\label{9-massless-particles}

The following particles are predicted to be \textbf{exactly massless} by symmetry:

{\def\LTcaptype{none} % do not increment counter
\begin{longtable}[]{@{}>{\RaggedRight\arraybackslash}p{0.480\textwidth}>{\RaggedRight\arraybackslash}p{0.480\textwidth}@{}}
\toprule\noalign{}
Particle & Reason \\
\midrule\noalign{}
\endhead
\bottomrule\noalign{}
\endlastfoot
Photon & Unbroken \(U(1)_{\text{em}}\) gauge invariance forbids a mass term \\
Gluons (8) & Unbroken \(SU(3)_c\) gauge invariance forbids a mass term \\
Graviton & Diffeomorphism invariance (emergent from entanglement equilibrium) forbids a mass term \\
\end{longtable}
}

These are \textbf{symmetry-protected zeros}, not accidental cancellations. The gauge invariances are derived from the gluing structure (§6.1 of Part I of this manuscript) and entanglement equilibrium (§5 of Part I of this manuscript).

\begin{center}\rule{0.5\linewidth}{0.5pt}\end{center}

\subsection{10. Complete Spectrum: OPH Outputs vs PDG}\label{10-complete-spectrum-oph-outputs-vs-pdg}

All outputs below use inputs \(P = 1.63094\) and \(\log(\dim\mathcal{H}) = 10^{122}\) only. PDG values are from the 2024/2025 edition. Each output is classified as \textbf{consistency check} (Tier~1, from calibration sector), \textbf{prediction} (Tier~2, structurally independent of calibration), or \textbf{structural} (Tier~3, parameter-free). See §1B for the full classification.

\subsubsection{10.1 Gauge Bosons and Higgs}\label{101-gauge-bosons-and-higgs}

{\def\LTcaptype{none} % do not increment counter
\begin{longtable}[]{@{}>{\RaggedRight\arraybackslash}p{0.160\textwidth}>{\RaggedLeft\arraybackslash}p{0.160\textwidth}>{\RaggedLeft\arraybackslash}p{0.160\textwidth}>{\RaggedLeft\arraybackslash}p{0.160\textwidth}>{\RaggedRight\arraybackslash}p{0.160\textwidth}>{\RaggedRight\arraybackslash}p{0.160\textwidth}@{}}
\toprule\noalign{}
Particle & OPH (GeV) & PDG (GeV) & Rel. Error & Tier & Error Source \\
\midrule\noalign{}
\endhead
\bottomrule\noalign{}
\endlastfoot
\(\gamma\) & 0 & 0 & exact & Structural & Symmetry-protected; no error expected \\
\(g\) (gluon) & 0 & 0 & exact & Structural & Symmetry-protected; no error expected \\
\(W\) boson & 80.386 & 80.377 ± 0.012 & +0.012\% & Consistency & Flows from calibrated \(\alpha_U\) and \(v\). Tree-level only \\
\(Z\) pole & 91.220 & 91.188 ± 0.002 & +0.035\% & Consistency & Flows from calibrated couplings. Simplified \(\Delta\rho\) \\
\(H\) (Higgs) & 126.48 & 125.20 ± 0.11 & +1.02\% & \textbf{Prediction} & Independent boundary condition \(\lambda = \beta_\lambda = 0\). 1-loop SM RGE only \\
\end{longtable}
}

\subsubsection{10.2 Charged Leptons}\label{102-charged-leptons}

{\def\LTcaptype{none} % do not increment counter
\begin{longtable}[]{@{}>{\RaggedRight\arraybackslash}p{0.160\textwidth}>{\RaggedLeft\arraybackslash}p{0.160\textwidth}>{\RaggedLeft\arraybackslash}p{0.160\textwidth}>{\RaggedLeft\arraybackslash}p{0.160\textwidth}>{\RaggedRight\arraybackslash}p{0.160\textwidth}>{\RaggedRight\arraybackslash}p{0.160\textwidth}@{}}
\toprule\noalign{}
Particle & OPH (GeV) & PDG (GeV) & Rel. Error & Tier & Error Source \\
\midrule\noalign{}
\endhead
\bottomrule\noalign{}
\endlastfoot
Electron & 5.109 × 10⁻⁴ & 5.110 × 10⁻⁴ & −0.023\% & \textbf{Prediction} & Koide structure (\(Q=2/3\), \(\delta=2/9\)) independent of calibration; residual from QED running \\
Muon & 0.10564 & 0.10566 & −0.022\% & \textbf{Prediction} & Same Koide structure; QED running corrections \textasciitilde0.02\%, matching the offset \\
Tau & 1.7766 & 1.7769 ± 0.09 & −0.020\% & \textbf{Prediction} & Same origin. Uniform sign shows common scale-factor effect from QED running \\
\end{longtable}
}

\subsubsection{10.3 Neutrinos}\label{103-neutrinos}

{\def\LTcaptype{none} % do not increment counter
\begin{longtable}[]{@{}>{\RaggedRight\arraybackslash}p{0.192\textwidth}>{\RaggedLeft\arraybackslash}p{0.192\textwidth}>{\RaggedLeft\arraybackslash}p{0.192\textwidth}>{\RaggedRight\arraybackslash}p{0.192\textwidth}>{\RaggedRight\arraybackslash}p{0.192\textwidth}@{}}
\toprule\noalign{}
Particle & OPH (GeV) & Experiment & Tier & Error Source \\
\midrule\noalign{}
\endhead
\bottomrule\noalign{}
\endlastfoot
\(\nu_e\) & 8.39 × 10⁻¹⁴ & \(< 8 \times 10^{-10}\) (KATRIN) & \textbf{Prediction} & Uses second input (\(\log\dim\mathcal{H}\)), independent of \(P\). Consistent with all bounds \\
\(\nu_\mu\) & 5.04 × 10⁻¹³ & \(< 8 \times 10^{-10}\) (KATRIN) & \textbf{Prediction} & Same; \(\Delta m_{21}^2\) consistent with solar oscillation scale in order of magnitude \\
\(\nu_\tau\) & 3.02 × 10⁻¹² & \(< 8 \times 10^{-10}\) (KATRIN) & \textbf{Prediction} & Anchored at \(\rho_\Lambda^{1/4}\); \(\Delta m_{32}^2\) consistent with atmospheric scale \\
\end{longtable}
}

All predictions are well below the current experimental upper bound of \(\sim 0.8\) eV \(\approx 8 \times 10^{-10}\) GeV. The predicted mass splittings are consistent with oscillation data in order of magnitude.

\subsubsection{10.4 Quarks}\label{104-quarks}

{\def\LTcaptype{none} % do not increment counter
\begin{longtable}[]{@{}>{\RaggedRight\arraybackslash}p{0.160\textwidth}>{\RaggedLeft\arraybackslash}p{0.160\textwidth}>{\RaggedLeft\arraybackslash}p{0.160\textwidth}>{\RaggedLeft\arraybackslash}p{0.160\textwidth}>{\RaggedRight\arraybackslash}p{0.160\textwidth}>{\RaggedRight\arraybackslash}p{0.160\textwidth}@{}}
\toprule\noalign{}
Particle & OPH (GeV) & PDG (GeV) & Rel. Error & Tier & Error Source \\
\midrule\noalign{}
\endhead
\bottomrule\noalign{}
\endlastfoot
up & 3.74 × 10⁻³ & 2.16 × 10⁻³ & +73\% & \textbf{Prediction} & \textbf{\(u\)-\(d\) degeneracy}: both get exponent \(n=6\). Subleading effects not included \\
down & 3.74 × 10⁻³ & 4.70 × 10⁻³ & −20\% & \textbf{Prediction} & Same degeneracy. Geometric mean \(\sqrt{m_u m_d} \approx 3.2\) MeV closer to 3.7 MeV \\
strange & 0.1346 & 0.0935 & +44\% & \textbf{Prediction} & No scheme matching; missing Clebsch factors and CKM rotation \\
charm & 0.808 & 1.273 & −37\% & \textbf{Prediction} & Missing RG running from \(v\) to \(m_c\) (\textasciitilde factor 1.3) and Clebsch coefficient \\
bottom & 4.847 & 4.183 & +16\% & \textbf{Prediction} & Running from \(v\) to \(m_b\) reduces mass by \textasciitilde15\%, closing most of the gap \\
top (texture) & 174.5 & 172.6 & +1.1\% & \textbf{Prediction} & Exponent \(n_t = 0\) (unsuppressed Yukawa); small error from \(v\) offset \\
top (crit. surf.) & 171.1 & 172.6 & −0.87\% & \textbf{Prediction} & Independent critical surface derivation; missing 2-loop RGE \\
\end{longtable}
}

\subsubsection{\texorpdfstring{10.5 Gauge Couplings at \(m_Z\)}{10.5 Gauge Couplings at m\_Z}}\label{105-gauge-couplings-at-m_z}

These are \textbf{calibration-sector consistency checks} (Tier~1). The gauge couplings at \(m_Z\) are the data from which \(P\) is extracted via the entropy-matching condition and self-consistent fixed-point procedure. Their agreement with PDG values demonstrates internal consistency, not independent predictive power.

{\def\LTcaptype{none} % do not increment counter
\begin{longtable}[]{@{}>{\RaggedRight\arraybackslash}p{0.18\textwidth}>{\RaggedLeft\arraybackslash}p{0.14\textwidth}>{\RaggedLeft\arraybackslash}p{0.18\textwidth}>{\RaggedLeft\arraybackslash}p{0.12\textwidth}>{\RaggedRight\arraybackslash}p{0.10\textwidth}>{\RaggedRight\arraybackslash}p{0.18\textwidth}@{}}
\toprule\noalign{}
Quantity & OPH & PDG & Rel. Error & Tier & Error Source \\
\midrule\noalign{}
\endhead
\bottomrule\noalign{}
\endlastfoot
\(\alpha_{\text{em}}^{-1}(m_Z)\) & 128.31 & 127.952 ± 0.009 & +0.28\% & Consistency & Flows from calibrated \(\alpha_U\) \\
\(\sin^2\theta_W(m_Z)\) & 0.2307 & 0.23122 ± 0.00004 & −0.21\% & Consistency & Most sensitive to threshold scale \\
\(\alpha_s(m_Z)\) & 0.1183 & 0.1179 ± 0.0009 & +0.37\% & Consistency & Within 0.5\(\sigma\) of PDG \\
\end{longtable}
}

\subsubsection{10.6 Derived Scales}\label{106-derived-scales}

{\def\LTcaptype{none} % do not increment counter
\begin{longtable}[]{@{}>{\RaggedRight\arraybackslash}p{0.14\textwidth}>{\RaggedLeft\arraybackslash}p{0.13\textwidth}>{\RaggedLeft\arraybackslash}p{0.13\textwidth}>{\RaggedLeft\arraybackslash}p{0.10\textwidth}>{\RaggedRight\arraybackslash}p{0.13\textwidth}>{\RaggedRight\arraybackslash}p{0.27\textwidth}@{}}
\toprule\noalign{}
Quantity & OPH Value & Reference & Agreement & Tier & Error Source \\
\midrule\noalign{}
\endhead
\bottomrule\noalign{}
\endlastfoot
\(v\) (Higgs VEV) & 246.77 GeV & 246.22 GeV & +0.22\% & Consistency & Flows from calibrated \(\alpha_U\) via transmutation formula \\
\(M_U\) & 2.47 × 10¹⁶ GeV & --- & --- & Consistency & No direct measurement; consistent with proton stability bounds \\
\(\alpha_U^{-1}\) & 24.32 & --- & --- & Consistency & Intermediate variable in calibration procedure \\
\(\Lambda_{\overline{\text{MS}}}^{(3)}\) & 0.322 GeV & \textasciitilde0.332 GeV & −3.0\% & Consistency & Propagated from \(\alpha_s\) (calibration sector); 0.4\% in \(\alpha_s\) \(\to\) \textasciitilde3\% in \(\Lambda\) \\
\end{longtable}
}

\subsubsection{10.7 Hadron Masses}\label{107-hadron-masses}

Hadron masses follow from \(m_X = C_X \cdot \Lambda_{\overline{\text{MS}}}^{(3)}\), where \(C_X\) is a dimensionless ratio computed via lattice QCD. The derivation chain is complete: \(P \to \alpha_s(m_Z) \to \Lambda^{(3)} \to m_X\). Current lattice estimates use small volumes and the quenched approximation, so systematic uncertainties are large (\textasciitilde10--20\%).

{\def\LTcaptype{none} % do not increment counter
\begin{longtable}[]{@{}>{\RaggedRight\arraybackslash}p{0.240\textwidth}>{\RaggedLeft\arraybackslash}p{0.240\textwidth}>{\RaggedRight\arraybackslash}p{0.240\textwidth}>{\RaggedRight\arraybackslash}p{0.240\textwidth}@{}}
\toprule\noalign{}
Particle & PDG (GeV) & Formula & Dominant uncertainty \\
\midrule\noalign{}
\endhead
\bottomrule\noalign{}
\endlastfoot
Proton & 0.93827 & \(C_p \cdot \Lambda^{(3)}\), \(C_p \approx 2.9\) & Quenching (\textasciitilde10\%), finite volume \\
Neutron & 0.93957 & \(\approx m_p\) in the isospin limit & Isospin splitting requires unquenched QCD+QED \\
\(\pi^\pm\) & 0.13957 & \(C_\pi \cdot \Lambda^{(3)}\) & Chiral extrapolation (\(m_\pi^2 \propto m_q\)) \\
\(\pi^0\) & 0.13498 & \(\approx m_{\pi^\pm}\) & EM splitting \\
\(K^\pm\) & 0.4937 & \(C_K \cdot \Lambda^{(3)}\) & Strange-quark valence mass \\
\(K^0\) & 0.4976 & \(\approx m_{K^\pm}\) & EM splitting \\
\(\Lambda\) baryon & 1.1157 & \(C_\Lambda \cdot \Lambda^{(3)}\) & Strange-quark baryon correlator \\
\end{longtable}
}

\begin{center}\rule{0.5\linewidth}{0.5pt}\end{center}

\subsection{11. Analysis of Discrepancies}\label{11-analysis-of-discrepancies}

\subsubsection{11.1 Sub-Permille Outputs (\textless{} 0.1\%)}\label{111-sub-permille-outputs--01}

\textbf{W boson} (+0.012\%), \textbf{Z boson} (+0.035\%), \textbf{electron} (−0.023\%), \textbf{muon} (−0.022\%), \textbf{tau} (−0.020\%):

These five outputs are at the sub-permille level, but their epistemic status differs. The gauge boson masses (\(W\), \(Z\)) are \textbf{consistency checks}: they flow directly from the calibrated gauge couplings and \(v\), so their agreement is expected by construction. The charged leptons (\(e\), \(\mu\), \(\tau\)) are \textbf{genuine predictions}: the Koide structure (\(Q = 2/3\), \(\delta = 2/9\)) is algebraically independent of the gauge-coupling calibration and has no continuous free parameter.

\subsubsection{11.2 Percent-Level Outputs (0.1\% -- 2\%)}\label{112-percent-level-outputs-01--2}

\textbf{Higgs} (+1.0\%), \textbf{top pole} (−0.87\%), \textbf{\(\alpha_s\)} (+0.37\%), \textbf{\(\sin^2\theta_W\)} (−0.21\%), \textbf{\(\alpha_{\text{em}}^{-1}\)} (+0.28\%):

The epistemic status differs within this group. The \textbf{gauge couplings} (\(\alpha_s\), \(\sin^2\theta_W\), \(\alpha_{\text{em}}^{-1}\)) are \textbf{consistency checks} --- they are the data used to calibrate \(P\), so their agreement demonstrates that the self-consistent fixed-point procedure works, not that the framework predicts them.

The \textbf{Higgs and top masses} are \textbf{genuine predictions}: they come from the critical surface constraint (\(\lambda(M_U) = 0\), \(\beta_\lambda(M_U) = 0\)) with 1-loop RG running --- an independent boundary condition not used in calibrating \(P\). The \textasciitilde1\% error is consistent with the expected size of 2-loop corrections and threshold effects.

\subsubsection{11.3 Qualitative Agreement Only: Quark Masses (16\% -- 73\%)}\label{113-qualitative-agreement-only-quark-masses-16--73}

The quark masses from the \(\varepsilon = 1/6\) texture show large deviations. The reasons are well-understood:

\begin{enumerate}
\def\labelenumi{\arabic{enumi}.}
\item
  \textbf{\(u\)-\(d\) degeneracy}: The texture assigns identical exponents \(n_u = n_d = 6\) to the up and down quarks, predicting \(m_u = m_d\). In reality, \(m_u/m_d \approx 0.46\). Breaking this degeneracy requires isospin-violating effects beyond the minimal \(\mathbb{Z}_6\) texture (e.g., subleading defect insertions or electromagnetic corrections).
\item
  \textbf{Scheme and scale mismatch}: The PDG quark masses are \(\overline{\text{MS}}\) running masses at scale \(\mu = 2\) GeV (for light quarks) or \(\mu = m_q\) (for heavy quarks). The OPH texture produces "Yukawa-scale" masses at \(\mu \sim v\) with no scheme matching applied. For the charm quark, this can easily account for a factor of \textasciitilde1.5.
\item
  \textbf{Order-one Clebsch-Gordan coefficients}: The texture \(y_f = c_f \cdot \varepsilon^{n_f}\) has residual coefficients \(c_f\) that are undetermined (expected to be \(\mathcal{O}(1)\), actually ranging from 0.6 to 2.2). These encode CKM rotation effects, RG running from \(M_U\) to \(m_Z\), and overlap matrix elements that the minimal integer-exponent ansatz does not resolve.
\item
  \textbf{Missing threshold corrections}: No QCD or electroweak threshold corrections are applied at flavor thresholds. These are typically 5--20\% effects for the lighter quarks.
\end{enumerate}

\textbf{The correct interpretation}: The texture correctly captures the \emph{hierarchy} (the fact that \(m_t/m_u \sim 10^5\) arises from integer exponents in base 6), and it does not aim for precision at the individual-mass level. The base-6 logarithms of all nine charged-fermion Yukawas land within \textasciitilde0.3 of integers , this is the core prediction. The order-one residuals are where scheme matching and subleading effects live.

\subsubsection{11.4 Hadron Masses}\label{114-hadron-masses}

The hadron mass predictions follow from the dimensionless ratios \(C_X = m_X/\Lambda^{(3)}\), computed via lattice QCD. The internal lattice code (\texttt{oph\_lattice\_su3\_quenched\_v5.py}) implements a quenched Wilson-gauge SU(3) lattice with Wilson valence quarks, gradient-flow scale setting, and Richardson continuum extrapolation , all without PDG inputs. The dominant systematic uncertainties are:

\begin{itemize}
\tightlist
\item
  \textbf{Quenching error} (\textasciitilde10\%): The gauge field is quenched (\(n_f = 0\)), while physical QCD has \(n_f = 2+1\) light dynamical flavors.
\item
  \textbf{Volume effects}: Lattice volumes (\(L = 2\)--\(6\)) give finite-volume corrections.
\item
  \textbf{Statistics}: Current runs use \(\mathcal{O}(10)\) configurations; scaling to \(\mathcal{O}(10^3)\) would reduce statistical errors.
\end{itemize}

These are standard lattice QCD systematics, not limitations of the OPH framework. The derivation chain \(P \to \alpha_s \to \Lambda^{(3)} \to m_X\) is complete; improving the lattice computation improves the numerical precision of the hadron predictions.

\subsubsection{11.5 Neutrinos: No Direct Comparison Available}\label{115-neutrinos-no-direct-comparison-available}

The predicted neutrino masses (\(\sim 0.08\)--\(3\) meV) are below current direct-detection sensitivity (KATRIN: \(m_\beta < 0.8\) eV) but in the right ballpark for cosmological and oscillation constraints. The predicted \(\sum m_\nu \approx 3.6\) meV is well below the cosmological bound \(\sum m_\nu < 0.12\) eV (Planck 2018).

\begin{center}\rule{0.5\linewidth}{0.5pt}\end{center}

\subsection{12. Reproduction Instructions}\label{12-reproduction-instructions}

All computations can be reproduced from the code in \href{../code/particles/}{\texttt{code/particles/}}.

\subsubsection{12.1 Prerequisites}\label{121-prerequisites}

\begin{verbatim}
pip install numpy
\end{verbatim}

No other dependencies are needed. The \texttt{pdg} and \texttt{pandas} packages are only required for \path{tools/fetch_pdg_data.py} (the PDG data fetcher, used for comparison only).

\subsubsection{12.2 Code Files}\label{122-code-files}

{\def\LTcaptype{none} % do not increment counter
\begin{longtable}[]{@{}>{\RaggedRight\arraybackslash}p{0.390\textwidth}>{\RaggedRight\arraybackslash}p{0.390\textwidth}>{\RaggedRight\arraybackslash}p{0.140\textwidth}@{}}
\toprule\noalign{}
File & Description & Stage \\
\midrule\noalign{}
\endhead
\bottomrule\noalign{}
\endlastfoot
\href{../code/particles/particle_masses_stage5.py}{\path{particle_masses_stage5.py}} & Core spectrum prediction: gauge closure, transmutation, critical surface, Z$_6$ texture & 1--5 \\
\href{../code/particles/oph_qcd.py}{\path{oph_qcd.py}} & 4-loop MSbar $\beta$-coefficients and $\Lambda$ extraction & 6 \\
\href{../code/particles/oph_lattice_su3_quenched_v5.py}{\path{oph_lattice_su3_quenched_v5.py}} & Quenched Wilson SU(3) lattice for hadron mass ratios & 6--7 \\
\href{../code/particles/oph_predict_compare.py}{\path{oph_predict_compare.py}} & Full predictor plus PDG comparison (main entry point) & All \\
\href{../code/particles/oph_no_cheat_audit.py}{\path{oph_no_cheat_audit.py}} & Static anti-leak audit of prediction code & Audit \\
\href{../code/particles/test_oph_predict_compare.py}{\path{test_oph_predict_compare.py}} & Smoke tests plus runtime no-cheat mutation test & Tests \\
\href{../code/particles/test_particle_masses_stage5.py}{\path{test_particle_masses_stage5.py}} & Regression tests for Stage 5 predictions & Tests \\
\end{longtable}
}

\subsubsection{12.3 Running the Full Prediction}\label{123-running-the-full-prediction}

\begin{Shaded}
\begin{Highlighting}[]
\BuiltInTok{cd}\NormalTok{ code/particles/}

\CommentTok{\# Full prediction with PDG comparison table}
\ExtensionTok{python3}\NormalTok{ oph\_predict\_compare.py }\AttributeTok{{-}{-}compare}

\CommentTok{\# JSON output for programmatic use}
\ExtensionTok{python3}\NormalTok{ oph\_predict\_compare.py }\AttributeTok{{-}{-}compare} \AttributeTok{{-}{-}json}

\CommentTok{\# With internal hadron computation (slow; tiny{-}lattice demo)}
\ExtensionTok{python3}\NormalTok{ oph\_predict\_compare.py }\AttributeTok{{-}{-}compare} \AttributeTok{{-}{-}with{-}hadrons} \AttributeTok{{-}{-}hadron{-}profile}\NormalTok{ demo}
\end{Highlighting}
\end{Shaded}

\subsubsection{12.4 No-Leakage Verification}\label{124-no-leakage-verification}

\begin{Shaded}
\begin{Highlighting}[]
\BuiltInTok{cd}\NormalTok{ code/particles/}

\CommentTok{\# Static audit: checks that build\_spectrum() does not reference PDG values}
\ExtensionTok{python3}\NormalTok{ oph\_no\_cheat\_audit.py}

\CommentTok{\# Runtime mutation test: scrambles PDG table and verifies predictions unchanged}
\ExtensionTok{python3}\NormalTok{ test\_oph\_predict\_compare.py}
\end{Highlighting}
\end{Shaded}

\subsubsection{12.5 Individual Stages}\label{125-individual-stages}

\begin{Shaded}
\begin{Highlighting}[]
\BuiltInTok{cd}\NormalTok{ code/particles/}

\CommentTok{\# Stage 5 spectrum (core charged sector)}
\ExtensionTok{python3}\NormalTok{ particle\_masses\_stage5.py}

\CommentTok{\# QCD scale extraction (standalone)}
\ExtensionTok{python3}\NormalTok{ oph\_qcd.py}
\end{Highlighting}
\end{Shaded}

\subsubsection{12.5 Fetching PDG Reference Data}\label{125-fetching-pdg-reference-data}

\begin{Shaded}
\begin{Highlighting}[]
\ExtensionTok{pip}\NormalTok{ install pdg pandas}
\ExtensionTok{python3}\NormalTok{ ../tools/fetch\_pdg\_data.py}
\CommentTok{\# Output: ../pdg\_data/particle\_masses.\{csv,json\}}
\end{Highlighting}
\end{Shaded}

\begin{center}\rule{0.5\linewidth}{0.5pt}\end{center}

\subsection{Summary}\label{summary}

Starting from one calibrated constant --- the pixel area \(P = 1.63094\) --- and one cosmological input --- the screen capacity \(\log(\dim\mathcal{H}) \sim 10^{122}\) --- the OPH derivation chain produces outputs in three epistemic tiers:

\textbf{Calibration-sector consistency checks} (agreement expected by construction):
\begin{itemize}\tightlist
\item Gauge couplings at \(m_Z\): \(\alpha_s\), \(\sin^2\theta_W\), \(\alpha_{\text{em}}^{-1}\) (0.2--0.4\%)
\item Electroweak boson masses: \(W\) (+0.012\%), \(Z\) (+0.035\%)
\item Higgs VEV \(v\) (+0.22\%)
\end{itemize}

\textbf{Genuine predictions} (structurally independent of calibration):
\begin{itemize}\tightlist
\item Charged leptons: \(e\), \(\mu\), \(\tau\) at \textless{} 0.025\% (from parameter-free Koide structure)
\item Higgs mass: 126.48 GeV vs 125.20 GeV (+1.0\%, from critical surface \(\lambda = \beta_\lambda = 0\))
\item Top quark: 171.1 GeV vs 172.6 GeV (−0.87\%, from critical surface)
\item 6 quark masses at correct order of magnitude with hierarchy reproduced, but 16--73\% individual errors
\item 3 neutrino masses consistent with all current bounds
\item Hadron masses from \(\Lambda_{\overline{\text{MS}}}^{(3)}\) via lattice QCD (10--20\% systematics)
\end{itemize}

\textbf{Parameter-free structural predictions} (from axioms alone):
\begin{itemize}\tightlist
\item Gauge group \(SU(3) \times SU(2) \times U(1)/\mathbb{Z}_6\), \(N_c = 3\), \(N_g = 3\)
\item 3 exactly massless particles (\(\gamma\), \(g\), graviton) from symmetry protection
\item Proton stability, \(\varepsilon = 1/6\), \(\delta = 2/9\), \(Q = 2/3\)
\end{itemize}

The derivation chain is:

\[
\boxed{
\begin{aligned}
P &\xrightarrow{\text{entropy matching}} \alpha_U, M_U
\xrightarrow{\text{transmutation}} v \\
&\xrightarrow{\text{RG + pixel}} \alpha_i(m_Z), m_Z, m_W
\xrightarrow{\text{critical surface}} m_H, m_t \\
&\xrightarrow{\mathbb{Z}_6\text{ texture}} m_f
\xrightarrow{\Lambda_{\overline{\text{MS}}}} m_{\text{hadrons}}
\end{aligned}
}
\]

\(P\) is the single free parameter, calibrated from the gauge coupling sector. The first arrow is calibration; subsequent arrows produce genuine predictions. The framework's claim is not ``we predict everything from nothing'' but: \textbf{from one calibrated constant and structural axioms, the framework determines the complete particle spectrum, with most outputs independent of the calibration.}
